\documentclass[a4paper,10pt]{article}
\hfuzz=10pt
\hbadness=10000

\usepackage{fontawesome5}
\usepackage{academicons}
\usepackage{graphicx}
\usepackage{url}
\usepackage{makecell}
\usepackage{parskip}
\usepackage{enumitem}
\setlist[itemize]{label=\raisebox{0.2ex}{\footnotesize\(\triangleright\)}}

\usepackage[usenames,dvipsnames]{xcolor}
\usepackage[scale=1.0, hmargin=0.75in, vmargin=0.5in]{geometry}
\usepackage{tabularx}
\usepackage{supertabular}
\usepackage{multicol}
\newcolumntype{C}{>{\centering\arraybackslash}X} 

\usepackage[unicode, draft=false]{hyperref}
\definecolor{ieeeblue0}{RGB}{0,67,147}
\definecolor{ieeeblue1}{RGB}{19,138,218}
\hypersetup{colorlinks,breaklinks,urlcolor=ieeeblue1,linkcolor=ieeeblue1}

\usepackage{tikz}
\newcommand{\sanf}[2]{{\textsf{\textcolor{#1}{#2}}}}

\usepackage{titlesec}
\titleformat{\section}
  {\sffamily\Large\color{ieeeblue0}}
  {}{0em}{}[\titlerule]
\titlespacing{\section}{0pt}{\baselineskip}{6pt}

\titleformat{\subsection}
  [runin]
  {\sffamily\normalsize\color{ieeeblue0}}
  {}{0em}{}[]
\titlespacing{\subsection}{0ex}{0ex}{0ex}

\newcommand{\badge}[2]{%
  \tikz[baseline=(n.base)]%
    \node (n) [%
      draw=#1!40,
      fill=#1!12,
      text=#1!50!black,
      rounded corners=4pt,
      inner xsep=4pt,
      inner ysep=0pt,
      minimum height=\baselineskip,
      text depth=0pt,
      font=\sffamily\normalsize
    ] {#2};%
}

% QR cmd
\newlength{\QRsize}
\setlength{\QRsize}{4cm}
\newcommand{\qrimg}[2]{%
  \href{#2}{\includegraphics[width=\QRsize,height=\QRsize]{#1}}%
}

\begin{document}

\pagestyle{empty} 

\newcommand{\icon}[1]{\raisebox{-0.25ex}{\includegraphics[height=1em]{#1}}}
\begin{minipage}[t]{0.4\linewidth}
  {\sffamily\bfseries\Huge\color{ieeeblue0}{Rui Luo}}
\end{minipage}%
\hfill
\begin{minipage}[t]{0.55\linewidth}
  \raggedleft\sffamily
  Email: \href{mailto:rui.1002@proton.me}{\textcolor{ieeeblue1}{rui.1002@proton.me}}\\
  Phone: \textcolor{ieeeblue1}{+86 13287816792}\\
  ShanghaiTech University, Shanghai, China
\end{minipage}
\vspace{1\baselineskip}
\\
\sanf{ieeeblue0}{\textbf{Summary:}} M.S. in Biomedical Engineering, ShanghaiTech University. Experienced in non-Cartesian pulse sequence design and reconstruction, with publications in IEEE TBME / ISMRM / SCMR / AAAI / JMRI, and open-source projects on PyPI / GitHub, invited as reviewer by IEEE TMI / ISMRM.

{\sanf{ieeeblue0}{\textbf{Skills:}}} Pulse Sequence Design, Reconstruction, Artifacts Correction, Simulation, Python (with C-API, CUDA acceleration), MATLAB (with MEX interface), C/C++, CUDA, Linux, Circuit Design, Embedded Programming, Desktop Development (Qt)

\section{Education}
\begin{tabularx}{\linewidth}{@{}l X@{}}	
2023 - 2026 & \sanf{ieeeblue0}{M.S. (Biomedical Engineering) at ShanghaiTech University} \hfill \normalsize \makecell[r]{(major GPA: 3.92/4.00)\\(overall GPA: 3.79/4.00)} \\
& \sanf{ieeeblue0}{Major Courses:} Algorithms \textbf{(A+)}, Principles of MRI \textbf{(A)}, Medical Image Processing \textbf{(A)}, \\
& MRI System Design \textbf{(A)}, Semiconductor Detectors and Readout Electronics \textbf{(A)} \\
\\
2018 - 2022 & \sanf{ieeeblue0}{B.S. (Automation) at Shandong University} \hfill \normalsize (overall GPA: 81.31/100) \\
& \sanf{ieeeblue0}{Awards:} $1\times$ National \textbf{1st} Prize, $1\times$ National \textbf{2nd} Prize, $1\times$ Provincial \textbf{1st} Prize.
\end{tabularx}

\section{Publications}
\subsection{[1] 2026, Real-Time Gradient Waveform Design for Arbitrary k-Space Trajectories}\;\badge{ieeeblue1}{Pulse Sequence}
\begin{itemize}
    \item \sanf{ieeeblue0}{Target:} Develop a general algorithm for calculating non-Cartesian gradient waveforms.
    \item \sanf{ieeeblue0}{Significance:} Improve the computational efficiency by \sanf{ieeeblue0}{1 order of magnitude}, achieving real-time gradient design. Improve the slew-rate accuracy by \sanf{ieeeblue0}{2 orders of magnitude}.
    \item \sanf{ieeeblue0}{Role:} \textbf{1st author}. Designed and implemented the proposed algorithm, and a trajectory library built on it.
    \item \sanf{ieeeblue0}{Status:} Published at \url{https://doi.org/10.1109/TBME.2026.3654117} on \textit{IEEE Trans. Biomed. Eng.}
\end{itemize}
\vspace{\baselineskip}
\subsection{[2] 2026, Sampling Density Compensation using Fast Fourier Deconvolution}\;\badge{ieeeblue1}{Reconstruction}
\begin{itemize}
    \item \sanf{ieeeblue0}{Target:} Develop a general algorithm for calculating density compensation functions.
    \item \sanf{ieeeblue0}{Significance:} Improve the computational efficiency by \sanf{ieeeblue0}{1-2 orders of magnitude}, saving tens of minutes.
    \item \sanf{ieeeblue0}{Role:} \textbf{1st author}. Designed and implemented the proposed algorithm.
    \item \sanf{ieeeblue0}{Status:} Published at \url{http://echo.ismrm.org/p/ISMRM2026/461-03-015} on \textit{ISMRM}-2026.
\end{itemize}
\vspace{\baselineskip}
\subsection{[3] 2026, Zero-shot INMR for Ultra-high Temporal Resolution Dynamic MRI}
\begin{itemize}
    \item \sanf{ieeeblue0}{Role:} \textbf{Coauthor}. Assistent in data acquisition using the algorithm in [1].
    \item \sanf{ieeeblue0}{Status:} Published at \url{https://doi.org/10.1609/aaai.v40i5.37395} on \textit{AAAI}-2026.
\end{itemize}
\vspace{\baselineskip}
\subsection{[4] 2026, Detection of Coronary Microvascular Dysfunction in Diabetic Mice Using Arterial Spin Labeling Cardiac MRI: A Multimodality Imaging Comparison}
\begin{itemize}
    \item \sanf{ieeeblue0}{Role:} \textbf{Coauthor}. Maintained an arterial perfusion marker (ASL) sequence on a Biospec 94/30 scanner.
    \item \sanf{ieeeblue0}{Status:} Published at \url{https://doi.org/10.1002/jmri.70304} on \textit{J. Magn. Reson. Imaging}.
\end{itemize}
\vspace{\baselineskip}
\subsection{[5] 2025, ASL-CMR detects CMD in diabetic mice: a multi-imaging modality comparison}
\begin{itemize}
    \item \sanf{ieeeblue0}{Role:} \textbf{Coauthor}. Maintained an arterial perfusion marker (ASL) sequence on a Biospec 94/30 scanner.
    \item \sanf{ieeeblue0}{Status:} Published at \url{https://archive.ismrm.org/2025/1028.html} on \textit{ISMRM}-2025.
\end{itemize}

% \clearpage

\section{Open-Source Projects}
\subsection{MRArbGrad}\;\badge{ieeeblue1}{Pulse Sequence}
\begin{itemize}
    \item \sanf{ieeeblue0}{Intro.:} The fastest non-Cartesian gradient solver at present; \sanf{ieeeblue0}{Real-time capability: gradients solve time is shorter than its duration;} Written in C++ backend with a Python wrapper; Compatible with all pulse sequence development platforms using C++ or Python; Includes a ready-to-use trajectory library; Manuscript Reference: [1].
    \item \sanf{ieeeblue0}{Link (PyPI): } \url{https://pypi.org/project/MRArbGrad}
    \item \sanf{ieeeblue0}{Link (GitHub): } \url{https://github.com/rui-luo1002/MRArbGrad}
    \item \sanf{ieeeblue0}{Role:} \textbf{Sole Developer}.
\end{itemize}
\vspace{\baselineskip}
\subsection{MRArbDcf}\;\badge{ieeeblue1}{Reconstruction}
\begin{itemize}
    \item \sanf{ieeeblue0}{Intro.:} The fastest density compensation function (DCF) solver at present; \sanf{ieeeblue0}{DCF solve time: 0.1 sec for 2D or 20 sec for 3D (CPU);} Supports both CPU and GPU platforms. Manuscript Reference: [2].
    \item \sanf{ieeeblue0}{Link (PyPI): } \url{https://pypi.org/project/MRArbDcf}
    \item \sanf{ieeeblue0}{Link (GitHub): } \url{https://github.com/rui-luo1002/MRArbDcf}
    \item \sanf{ieeeblue0}{Role:} \textbf{Sole Developer}.
\end{itemize}
\vspace{\baselineskip}
\subsection{TorchFinufft-Ryan}\;\badge{ieeeblue1}{Optimization}
\begin{itemize}
    \item \sanf{ieeeblue0}{Intro.:} NUFFT module for PyTorch and a Toeplitz MSE loss module; Useful in optimization-based or deep-learning-based reconstruction projects; \sanf{ieeeblue0}{In 256\textsuperscript{2} iterative reconstruction: 2 ms per step using NUFFT or 3 ms per step using Toeplitz on an RTX3090}.
    \item \sanf{ieeeblue0}{Link (PyPI): } \url{https://pypi.org/project/torchfinufft-ryan}
    \item \sanf{ieeeblue0}{Link (GitHub): } \url{https://github.com/rui-luo1002/torchfinufft}
    \item \sanf{ieeeblue0}{Role:} \textbf{Sole Developer}.
\end{itemize}
\vspace{\baselineskip}
\subsection{MRPhantom}\;\badge{ieeeblue1}{Simulation}
\begin{itemize}
    \item \sanf{ieeeblue0}{Intro.:} Volumetric dynamic MRI Phantom of a Slime with respiratory and cardiac motion, and M0, phase, T1, T2, B0 maps, and also coil sensitivity maps; Boosted by a parallel C-API Backend; \sanf{ieeeblue0}{Support memory-efficient mode: phantoms are generated in real time according to motion parameter to avoid explicit 3D+T construction (137 ms per 256\textsuperscript{3} volume).}
    \item \sanf{ieeeblue0}{Link (PyPI): } \url{https://pypi.org/project/MRPhantom}
    \item \sanf{ieeeblue0}{Link (GitHub): } \url{https://github.com/rui-luo1002/MRPhantom}
    \item \sanf{ieeeblue0}{Role:} \textbf{Sole Developer}.
\end{itemize}
\vspace{\baselineskip}

\section{Awards}
\subsection{National 1st prize in Smart Car Competition, Acoustic Group, 2020}\;\badge{ieeeblue1}{Signal Processing}
\begin{itemize}
    \item \sanf{ieeeblue0}{Target:} Design a model car to hit certain beacons playing acoustic chirp signals.
    \item \sanf{ieeeblue0}{Selectivity:} 1 team per category per school.
    \item \sanf{ieeeblue0}{Role:} \textbf{Team member}. Developed the fast cross-correlation algorithm for acoustic navigation on a microcontroller, and also a secondary optical navigation algorithm.
\end{itemize}
\vspace{\baselineskip}
\subsection{National 2nd prize in Smart Car Competition, Delivery Group, 2021}\;\badge{ieeeblue1}{Circuit Design}
\begin{itemize}
    \item \sanf{ieeeblue0}{Target:} Design a model car with LiDAR and Linux microcomputer to deliver packages in a maze.
    \item \sanf{ieeeblue0}{Selectivity:} 1 team per category per school.
    \item \sanf{ieeeblue0}{Role:} \textbf{Team leader}. Full-stack experience including circuit design and assembly, robot operating system debugging, optical navigation and PID algorithm development, TCP host application development, hardware driver development.
\end{itemize}

% \vspace{\baselineskip}
% \clearpage

\subsection{Provincial 1st prize in Smart Car Competition, Optical Group, 2019}\;\badge{ieeeblue1}{Embedded Programming}
\begin{itemize}
    \item \sanf{ieeeblue0}{Target:} Design a model car to trace the racing track using a digital camera.
    \item \sanf{ieeeblue0}{Selectivity:} 2 teams per category per school.
    \item \sanf{ieeeblue0}{Role:} \textbf{Team member}. Developed the optical navigation algorithm using Otsu's threshold algorithm, edge detection algorithm and PID control algorithm.
\end{itemize}
\sanf{ieeeblue0}{Other Awards:} Two additional national prizes and three university-level prizes for which I was invited to provide technical support.

\section{Experience}
\subsection{2025 September, Reviewer, \textit{IEEE Transactions on Medical Imaging}}
\begin{itemize}
    \item \sanf{ieeeblue0}{Role:} \textbf{Reviewer}.
    \item \sanf{ieeeblue0}{Highlights:} Reviewed one manuscript.
\end{itemize}
\vspace{\baselineskip}
\subsection{2025 November, Reviewer, \textit{International Society of Magnetic Resonance in Medicine}}
\begin{itemize}
    \item \sanf{ieeeblue0}{Role:} \textbf{Reviewer}.
    \item \sanf{ieeeblue0}{Highlights:} Reviewed 20 abstracts.
\end{itemize}
\vspace{\baselineskip}
\subsection{2024 Fall, Teaching Assistant, \textit{Principles of Magnetic Resonance Imaging}}
\begin{itemize}
    \item \sanf{ieeeblue0}{Role:} \textbf{Teaching Assistant}.
    \item \sanf{ieeeblue0}{Highlights:} Mentored 33 students; designed the final project involving sampling and reconstruction simulation; delivered a review lecture highly rated by students.
\end{itemize}

\section{Reference}
\begin{center}
\setlength{\tabcolsep}{10pt}
\begin{tabular}{ccc}
\qrimg{../qr/scholar.png}{https://scholar.google.com/citations?user=llTqHTUAAAAJ\&hl=en} &
\qrimg{../qr/pypi.png}{https://pypi.org/user/rui1002/} &
\qrimg{../qr/github.png}{https://github.com/rui-luo1002}
\end{tabular}
\end{center}

\end{document}
